
\documentclass{article}

\usepackage[a4paper, left=1.5cm, right=1.5cm, top=2cm, bottom=2cm]{geometry}
\usepackage{../../../../components/components}

\usepackage{fancyhdr}

% Configuration des en-têtes et pieds de page
\pagestyle{fancy}
\fancyhf{} % reset tout

\fancyhead[L]{Préparation CAPES}
\fancyhead[C]{Analyse}
\fancyhead[R]{2026-2027}

\fancyfoot[L]{Ewen Rodrigues de Oliveira}
\fancyfoot[R]{\thepage}

\begin{document}

\docTitle{Chapitre 4 : Interpolation de Lagrange}



\definition{
On considère \(n+1\) points distincts \((x_0,y_0), \dots, (x_n,y_n)\).
L’interpolation de Lagrange consiste à construire un polynôme \(P\) tel que :
\[
P(x_i)=y_i \quad \forall i \in \{0,\dots,n\}.
\]
}

\method{Construction du polynôme de Lagrange}{
On construit :
\[
P(x)=\sum_{i=0}^{n} y_i L_i(x)
\]
où les polynômes de base sont définis par :
\[
L_i(x)=\prod_{\substack{j=0 \\ j\neq i}}^{n} \frac{x-x_j}{x_i-x_j}.
\]
}

\remark{
Chaque \(L_i(x)\) vérifie :
\[
L_i(x_i)=1 \quad \text{et} \quad L_i(x_j)=0 \ (j\neq i).
\]
Ainsi, chaque point \((x_i,y_i)\) “active” uniquement son propre terme.
}

\example{
Soient les points \((1,2)\) et \((3,4)\).

On a :
\[
L_0(x)=\frac{x-3}{1-3}, \quad L_1(x)=\frac{x-1}{3-1}.
\]

Donc :
\[
P(x)=2L_0(x)+4L_1(x).
\]
Après simplification :
\[
P(x)=x+1.
\]
}

\attention{
Plus le nombre de points augmente, plus le polynôme devient de degré élevé et peut osciller fortement (phénomène de Runge).
}

\method{Procédure de calcul}{
\begin{enumerate}
    \item Lister les points \((x_i,y_i)\).
    \item Construire chaque \(L_i(x)\).
    \item Multiplier chaque \(L_i(x)\) par \(y_i\).
    \item Additionner tous les termes.
    \item Simplifier.
\end{enumerate}
}
\newpage
\training{
Calculer le polynôme de Lagrange passant par les points :
\[
(0,1), \ (1,2), \ (2,5).
\]
}
\noindent \carreaux{10}

\training{
Soient les points \((-1,1), (0,0), (1,1)\).
\begin{enumerate}
    \item Construire les polynômes \(L_i(x)\). \\
    \item En déduire \(P(x)\). \\
    \item Vérifier que \(P(x)\) est un polynôme du second degré.
\end{enumerate}
}
\noindent \carreaux{10}

\problem{
On considère les points \((1,1), (2,4), (3,9), (4,16)\).
Sans calcul complet, que peut-on conjecturer sur le degré du polynôme interpolateur ?
}

\end{document}